% Page budget: exactly 1.5 pages; at least two figures; no bullets.
\clearpage
\section{Background}

\subsection{LLM Serving and Request Scheduling}
% [作者手写区] 要点: 解释 continuous batching、PagedAttention、FCFS 与长度感知调度的必要概念；只保留理解本系统所需定义。Bib keys: vllm2023, fuLTR2024, ssjf2024.

\ReportFigure{fig1}{System architecture and request path.}{\texttt{scripts/plot\_final\_report\_figures.py --figure fig1}}{fig:system-architecture}

\subsection{Prediction and Reliability Signals}
% [作者手写区] 要点: 区分 absolute length、relative rank、uncertainty/OOD signal；说明 score 越大预测越长且 live scheduler 只消费 gateway verdict。Bib keys: fuLTR2024, pars2025, tie2026, egtp2026, beyondPrediction2026.

\subsection{Gateway and Engine Boundary}
% [作者手写区] 要点: 对照 Fig.1 说明 VeloxMesh 白名单、decision contract、vllm_xargs 和 fail-open；只写已经实现且有测试的能力。版本证据在 Tab.1；gateway base 4b4b5ad，LTR branch pin d49d79d。Bib keys: vllm2023.

\ReportFigure{fig2}{Training-artifact and serving-artifact lineage.}{\texttt{scripts/plot\_final\_report\_figures.py --figure fig2}}{fig:data-lineage}

\subsection{Agentic Context and KV State}
% [作者手写区] 要点: 说明 history、tool schema、tool result 和 KV/session state 为什么改变预测与调度边界；KV 协同本期若无实验只作为背景/未来工作。Bib keys: agentix2026, infercept2024, toolace2024, bfcl, toolathlon2025.
